\newcommand{\Dupples}{\texttt{Dupples}}
\newcommand{\Nuzzles}{\texttt{Nuzzles}}
\newcommand{\Prittles}{\texttt{Prittles}}
\newcommand{\Skipples}{\texttt{Skipples}}
\newcommand{\Crottles}{\texttt{Crottles}}

\section{Exercise 1}

    \subsection{Approach}

        I started out by modeling the products first, so I made an enum and
        gave it properties such as name, weight, and required amount. Then
        all I had to do was iterate over them and creating variables using
        the weight and the required amount for each product.

        Next, I started with the trucks. Trucks have a maximum weight they can
        carry, plus an amount for each product they're carrying. So I added
        those two fields and made the variables for: the total weight in the
        truck, the amount of product $p_1$, the amount for product $p_2$, ...
        
        To finish the trucks off I just had to assert that the weight is
        between 0 and the max capacity, and that it is equal to the sum of
        the weights of the products being transported. And finally the amount
        of products being transported had to be less the 10.

        The last part was just asserting a couple more global constraints. Like
        how the total amount of a product being transported is the same as the
        sum of what the trucks are transporting.

    \subsection{Implementation}

        \subsubsection{constants}
            \begin{itemize}
                \item
                    The amount of trucks as
                    $\texttt{amount\_of\_trucks} = 6$
                \item 
                    The maximum load capacity per truck in kg
                    $\texttt{max\_weight\_truck} = 8000$
                \item
                    The maximum amount of products in a truck
                    $\texttt{max\_amount\_truck} = 10$
            \end{itemize}

        \subsubsection{Variables}
            \begin{itemize}
                \item
                    $\alpha_{p,t} \in \mathbb{Z}$ representing the total amount
                    of a product $p$ being transported in a truck $t$.
                \item
                    $\beta_{p} \in \mathbb{Z}$ representing the total amount of
                    a product $p$ being transported.
                \item
                    $\gamma_{t} \in \mathbb{Z}$ representing the total weight a
                    truck $t$ is transporting.
            \end{itemize}

        \subsubsection{Sets}
            \begin{itemize}
                \item 
                    The set of products: $P = \big\{
                        \Dupples,
                        \Nuzzles,
                        \Prittles,
                        \Skipples,
                        \Crottles
                    \big\}$.
                \item 
                    The set of products that have a required amount: 
                    $P_r = P \setminus \{ \Dupples \}$
                \item
                    The set of trucks: $T = \big\{
                        1, 2, ..., \texttt{amount\_of\_trucks}
                    \big\}$
                \item
                    The set of trucks that can transport \Skipples:
                    $T_s = \big\{ 1, 2 \big\}$
            \end{itemize}

        \subsubsection{Functions}        
            \begin{itemize}
                \item 
                    A function to get the required amount of a product:
                    \[
                        r : P_r \to \mathbb{N},\qquad r(p) = \begin{cases}
                            6  & \text{if } p = Nuzzles  \\
                            12 & \text{if } p = Prittles \\
                            15 & \text{if } p = Skipples \\
                            8  & \text{if } p = Crottles
                        \end{cases}
                    \]
                \item 
                    A function to get the weight per product:
                    \[
                        w : P \to \mathbb{N},\qquad w(p) = \begin{cases}
                            800  & \text{if } p = Nuzzles  \\
                            405  & \text{if } p = Prittles \\
                            500  & \text{if } p = Skipples \\
                            2500 & \text{if } p = Crottles \\
                            600  & \text{if } p = Dupples \\
                        \end{cases}
                    \]
            \end{itemize}
        
        \subsubsection{Assertions}
            For the total amount of product $p$ being transported I created the
            following formulas:

            \begin{itemize}
                \item 
                    Asserting that the amount of products being transported are
                    non-negative:
                    \[
                        \bigg(
                            \bigwedge_{(p, t) \in P \times T} 
                                \alpha_{p,t} \ge 0
                        \bigg) \wedge \bigg(
                            \bigwedge_{p \in P} \beta_p \ge 0
                        \bigg)
                    \]
                \item
                    Asserting that the required amount is transported:
                    \[
                        \bigwedge_{p \in P_r} \beta_{p} = r(p)
                    \]
                \item
                    Asserting that the weight of a truck is the sum of its
                    products, and is in a valid range:
                    \[
                        \bigwedge_{t \in T} 0 \leq \gamma_{t} = \bigg(
                            \sum_{p \in P} \Big(
                                \alpha_{p,t}\cdot w(p)
                            \Big)
                        \bigg) \leq \texttt{max\_weight\_truck}
                    \]
                \item
                    asserting a truck $t$ transports at most 10 items:
                    \[
                        \bigwedge_{t \in T} \bigg(
                            \sum_{p \in P} \alpha_{p,t}
                        \bigg) \leq 10
                    \]
                \item
                    asserting that only some trucks can transport
                    \texttt{Skipples}:
                    \[
                        \bigwedge_{t \in T} \bigg(
                            \big(
                                t \not\in T_s 
                            \big) \implies \big(
                                \alpha_{p,t} = 0
                            \big)
                        \bigg) \qquad \text{ where $p$ = \Skipples}
                    \]
                \item
                    asserting the total amount of a product $p$ being
                    transported is the sum of what the trucks are transporting:
                    \[
                        \bigwedge_{p \in P} \bigg(
                            \beta_p = \Big( \sum_{t \in T} \alpha_{p,t}
                            \Big)
                        \bigg)
                    \]
                \item 
                    asserting that \Prittles\text{ } are spread over at least
                    5 trucks:
                    \[
                        \sum_{t \in T} \bigg(
                            \operatorname{ite}\Big(
                                \big(\alpha_{p,t} > 0\big),\ 1,\ 0
                            \Big)
                        \bigg) \geq 5 \qquad \text{ where p = \Prittles}
                    \]
            \end{itemize}

            \noindent The goal is of course to maximise the amount of \Dupples:
            
            \[
                \text{maximize } \beta_p \text{ where $p = \Dupples$}
            \]

            \noindent For question b there was one more asserting we have to
            make. Namely, \Crottles cannot be transported without at least 2
            \Dupples\text{ } in the same truck:

            \[
                \bigwedge_{t \in T} \bigg(
                    \Big(
                        \alpha_{p_1,t} > 0
                    \Big) \implies \Big(
                        \alpha_{p_2,t} \geq 2
                    \Big)
                \bigg) \qquad 
                    \text{ where $p_1 = \Crottles$}
                    \text{ and $p_2 = \Dupples$}
            \]

            \noindent I did leave out all mentions of truck speed. I did this
            on purpose as the driving speed does not affect the amount of
            \Dupples\text{ } I can deliver and that would just slow the solver
            down, and make it harder to implement.

    \subsection{Results}

        The results were that both sub-questions were satisfiable. I've created
        tables to help visualize the results.

        \subsubsection{Visualisation Exercise 1a}

            \begin{table}[ht]
                \centering
                \caption{Distribution of products over the trucks for 1a.}
                \label{tab:exercise1a}
                \begin{tabular}{c|ccccc|c}
                    \textbf{Truck} & \textbf{Nuzzles} & \textbf{Prittles} & 
                    \textbf{Skipples} & \textbf{Crottles} & \textbf{Dupples} &
                    \textbf{Weight (kg)} \\
                    \hline
                    %%% T | N | P | S | C | D | W
                        1 & 0 & 1 & 6 & 1 & 2 & 7105 \\
                        2 & 0 & 0 & 9 & 1 & 0 & 7000 \\
                        3 & 2 & 2 & 0 & 1 & 5 & 7910 \\
                        4 & 0 & 1 & 0 & 2 & 4 & 7805 \\
                        5 & 4 & 4 & 0 & 1 & 1 & 7920 \\
                        6 & 0 & 4 & 0 & 2 & 2 & 7820 \\
                    \hline
                    \textbf{Total} & \textbf{6} & \textbf{12} & \textbf{15} &
                    \textbf{8} & \textbf{14} & \\
                \end{tabular}
            \end{table}

        \subsubsection{Visualisation Exercise 1b}

            \begin{table}[ht]
                \centering
                \caption{Distribution of products over the trucks for 1b.}
                \label{tab:exercise1b}
                \begin{tabular}{c|ccccc|c}
                    \textbf{Truck} & \textbf{Nuzzles} & \textbf{Prittles} & 
                    \textbf{Skipples} & \textbf{Crottles} & \textbf{Dupples} &
                    \textbf{Weight (kg)} \\
                    \hline
                    %%% T | N | P | S | C | D | W
                        1 & 0 & 1 & 6 & 1 & 2 & 7105 \\
                        2 & 1 & 0 & 9 & 0 & 0 & 5300 \\
                        3 & 0 & 4 & 0 & 2 & 2 & 7820 \\
                        4 & 3 & 3 & 0 & 1 & 3 & 7915 \\
                        5 & 1 & 2 & 0 & 2 & 2 & 7810 \\
                        6 & 1 & 2 & 0 & 2 & 2 & 7810 \\
                    \hline
                    \textbf{Total} & \textbf{6} & \textbf{12} & \textbf{15} &
                    \textbf{8} & \textbf{11} & \\
                \end{tabular}
            \end{table}
